\documentclass[11pt]{article}
% Shared preamble for "The Kodamai Container Program: A Complete Technical Record"
% Self-contained; compiles with pdflatex.
\usepackage[utf8]{inputenc}
\usepackage[T1]{fontenc}
\usepackage{lmodern}
\usepackage{amsmath,amssymb,amsthm}
\usepackage{stmaryrd}   % \llbracket \rrbracket
\usepackage{mathrsfs}
\usepackage[margin=1in]{geometry}
\usepackage{enumitem}
\usepackage{microtype}
\usepackage{newunicodechar}
\usepackage[colorlinks=true,linkcolor=blue,citecolor=blue,urlcolor=blue]{hyperref}

% ---- Theorem environments -------------------------------------------------
\theoremstyle{plain}
\newtheorem{theorem}{Theorem}[section]
\newtheorem{proposition}[theorem]{Proposition}
\newtheorem{corollary}[theorem]{Corollary}
\newtheorem{lemma}[theorem]{Lemma}
\theoremstyle{definition}
\newtheorem{definition}[theorem]{Definition}
\newtheorem{result}[theorem]{Result}
\theoremstyle{remark}
\newtheorem*{remark}{Remark}
\newtheorem*{flag}{Flag}

% ---- Grade / registry helper macros ---------------------------------------
% Trust grades are copied verbatim from the master outline; never inflated.
\newcommand{\grade}[1]{\textbf{[#1]}}
\newcommand{\node}[1]{\texttt{#1}}

% ---- Unicode -> LaTeX mapping (so the outline prose transcribes verbatim) --
% Superscripts / subscripts
\newunicodechar{¹}{\ensuremath{^{1}}}
\newunicodechar{²}{\ensuremath{^{2}}}
\newunicodechar{³}{\ensuremath{^{3}}}
\newunicodechar{⁰}{\ensuremath{^{0}}}
\newunicodechar{⁺}{\ensuremath{^{+}}}
\newunicodechar{₁}{\ensuremath{_{1}}}
\newunicodechar{₂}{\ensuremath{_{2}}}
\newunicodechar{₃}{\ensuremath{_{3}}}
\newunicodechar{₄}{\ensuremath{_{4}}}
\newunicodechar{₊}{\ensuremath{_{+}}}
% Greek
\newunicodechar{Γ}{\ensuremath{\Gamma}}
\newunicodechar{Δ}{\ensuremath{\Delta}}
\newunicodechar{Π}{\ensuremath{\Pi}}
\newunicodechar{Σ}{\ensuremath{\Sigma}}
\newunicodechar{Ψ}{\ensuremath{\Psi}}
\newunicodechar{δ}{\ensuremath{\delta}}
\newunicodechar{ε}{\ensuremath{\varepsilon}}
\newunicodechar{θ}{\ensuremath{\theta}}
\newunicodechar{λ}{\ensuremath{\lambda}}
\newunicodechar{ν}{\ensuremath{\nu}}
\newunicodechar{φ}{\ensuremath{\varphi}}
\newunicodechar{ω}{\ensuremath{\omega}}
% Blackboard / script / fraktur
\newunicodechar{ℕ}{\ensuremath{\mathbb{N}}}
\newunicodechar{ℤ}{\ensuremath{\mathbb{Z}}}
\newunicodechar{ℰ}{\ensuremath{\mathcal{E}}}
\newunicodechar{𝒟}{\ensuremath{\mathcal{D}}}
\newunicodechar{𝒯}{\ensuremath{\mathcal{T}}}
\newunicodechar{𝔤}{\ensuremath{\mathfrak{g}}}
\newunicodechar{𝔽}{\ensuremath{\mathbb{F}}}
% Arrows and relations
\newunicodechar{→}{\ensuremath{\to}}
\newunicodechar{↔}{\ensuremath{\leftrightarrow}}
\newunicodechar{⟹}{\ensuremath{\implies}}
\newunicodechar{⟺}{\ensuremath{\iff}}
\newunicodechar{≃}{\ensuremath{\simeq}}
\newunicodechar{≅}{\ensuremath{\cong}}
\newunicodechar{≠}{\ensuremath{\neq}}
\newunicodechar{≡}{\ensuremath{\equiv}}
\newunicodechar{≤}{\ensuremath{\leq}}
\newunicodechar{≥}{\ensuremath{\geq}}
\newunicodechar{⊇}{\ensuremath{\supseteq}}
\newunicodechar{⊊}{\ensuremath{\subsetneq}}
\newunicodechar{∈}{\ensuremath{\in}}
% Operators
\newunicodechar{×}{\ensuremath{\times}}
\newunicodechar{∂}{\ensuremath{\partial}}
\newunicodechar{−}{\ensuremath{-}}
\newunicodechar{∘}{\ensuremath{\circ}}
\newunicodechar{∞}{\ensuremath{\infty}}
\newunicodechar{∧}{\ensuremath{\wedge}}
\newunicodechar{∨}{\ensuremath{\vee}}
\newunicodechar{⊕}{\ensuremath{\oplus}}
\newunicodechar{⊗}{\ensuremath{\otimes}}
\newunicodechar{⊙}{\ensuremath{\odot}}
\newunicodechar{⋆}{\ensuremath{\star}}
\newunicodechar{⋈}{\ensuremath{\bowtie}}
\newunicodechar{⋉}{\ensuremath{\ltimes}}
\newunicodechar{⋊}{\ensuremath{\rtimes}}
\newunicodechar{▷}{\ensuremath{\triangleright}}
\newunicodechar{◁}{\ensuremath{\triangleleft}}
\newunicodechar{⟦}{\ensuremath{\llbracket}}
\newunicodechar{⟧}{\ensuremath{\rrbracket}}
% Punctuation
\newunicodechar{–}{--}
\newunicodechar{—}{---}
\newunicodechar{§}{\S}


\title{The Kodamai Container Program: A Complete Technical Record\\[4pt]
\large Part IV --- Applications}
\author{MacBeth}
\date{2026-10-10}

\begin{document}
\maketitle
\begin{center}\itshape Prepared for Neil Ghani.\end{center}

\begin{abstract}
Part~IV records the applications of the container program: the same $\mathrm{DCont}\simeq
\mathrm{Cat}$ $+$ ZS-weld machinery models real compositional systems where correctness has
economic consequences. Each application is a directed container; composition is a ZS product;
failure modes are obstruction classes. Every result carries its trust grade verbatim and its
registry node; external/collaborator work and genuine write gaps are flagged as such.
\end{abstract}

\paragraph{Grade legend.}
$\grade{proved} > \grade{lean-verified} / \grade{computed}$. Grades are copied verbatim and never
inflated.

\paragraph{Thesis of Part IV.}
The same $\mathrm{DCont}\simeq\mathrm{Cat}$ $+$ ZS-weld machinery models real compositional systems
where correctness has economic consequences. Each application is a directed container; composition
is a ZS product; failure modes are obstruction classes.

\section{Supply chains --- inventory at $H^1$ and the gerbe at $H^2$}
\emph{Precis.} A supply chain is a small category / directed container; composing flows sharing a
warehouse is a ZS product. Inventory consistency is a \textbf{degree-one} $H^1$ datum, INDEPENDENT
of the weld class $[ω]\in H^2$. A finer relabelling-group obstruction lives at $H^2$ as a gerbe.

\begin{theorem}[Supply chain as ZS weld, \grade{computed}]
A supply chain is a directed container; composing two flows $=$ the ZS product $C\bowtie D$; the
$[ω]\in H^2$ obstruction applies. Node \node{supply-chain-zs}.
\end{theorem}

\begin{theorem}[Tier-2 inventory gerbe, \grade{proved}]
When local inventory charts agree only up to an iso of stock-states (relabelling group
$K=\mathrm{Aut}(F)$), the gluing obstruction is a gerbe class
$[𝔤]\in H^2(N(C);Z(\mathrm{Aut}\,F))$; it is rectifiable $\iff[𝔤]=0$, independent of the weld, and
nonvacuous $\iff Z(\mathrm{Aut}\,F)\neq1$. Node \node{supply-chain-inventory-gerbe-h2}.
\end{theorem}

\begin{theorem}[Inventory ${[θ_R]\in H^1}$, independent of the weld]
Inventory consistency is a degree-one datum $[θ_R]\in H^1$, INDEPENDENT of the weld
$[ω]\in H^2$ (proof file \texttt{supply-chain-inventory-degree-one.tex}; appears as the
$N_2$/$H^1$ rung of the nerve tower, node \node{composition-nerve-level}).
\end{theorem}
\begin{flag}
This $H^1$ result has \textbf{no standalone registry node} of its own --- it is carried inside
\node{composition-nerve-level}. It must not be cited as an independent \grade{proved} node without
opening its own registry entry.
\end{flag}

\section{Blockchain / smart-contract re-entrancy}
\emph{Precis.} Smart contracts as coalgebras; a re-entrancy bug is a failed distributive law / an
unprotected shared position in a ZS weld. The obstruction is machine-checkable.

\begin{theorem}[Re-entrancy as failed guard, \grade{proved}; class \grade{lean-verified}]
Re-entrancy $=$ a failed guard in the orchestration ZS weld (node \node{orchestration-zs}); the
$H^2$ class $[ω]=ε$ is machine-checked (\texttt{reentrancy-machine-checked-obstruction.tex} /
node \node{fibre-h2-z2-engine}).
\end{theorem}
\begin{flag}
Per Part~II, state $[ω]$ as a \emph{splitting} obstruction of an associative weld, not as
``non-associativity.''
\end{flag}

\begin{remark}[Security-audit framing]
State-divergence bugs are failed distributive laws (connection
\node{security-audit-bugs-are-failed-distributive-laws}; paper
\texttt{state-divergence-distributive-laws.tex}).
\end{remark}

\section{GA composition / linear attention (adjacent)}
\emph{Precis.} Island-model GA migration topology and linear-attention composition are instances of
the same compositional algebra in $\mathrm{Mat}(\mathrm{Vec})$ / Poly.

\begin{theorem}[Linear attention as $\mathrm{Mat}(\mathrm{Vec})$ composition, \grade{proved}]
Linear attention $=$ composition in the Vec-matrix bicategory $\mathrm{Mat}(\mathrm{Vec})$,
$(P\oplus Q)(a,c)=\oplus_b P(a,b)\otimes Q(b,c)$. Node \node{vec-attention-composition}. (The
``stack $=$ one matrix'' reading is REFUTED --- MEMORY.)
\end{theorem}

\begin{remark}[External / empirical]
The GA diversity-dynamics empirical result (Langer--Vega, GECCO 2026) and ga-containers (ACT 2026)
are external/empirical collaborator publications --- cite as published collaborator work, not a
MacBeth registry node.
\end{remark}

\section{Agent orchestration and tax computation}
\emph{Precis.} Tool-calling protocols and meta-agent orchestration are directed containers; tax
computation (Catala-style) is a scope/supersession DAG with a comonadic event-sourced interface.

\begin{remark}[Tool-calling / harness]
Tool-calling protocols are directed containers (shipped note
\texttt{tool-calling-protocols-directed-containers.tex}; grounded in \node{orchestration-zs},
\grade{proved}). The harness-is-a-worker / MCP formalisation (paper \texttt{harness-is-a-worker.tex};
connection \node{harness-worker-is-the-unscooped-mcp-formalization}).
\end{remark}

\begin{flag}[Tax computation --- write gap]
Named in SEED Path 5 but has \textbf{NO registry node and no MacBeth write-up}. Either produce the
ontology-DAG / supersession-DAG / event-sourced-comonad model or mark it explicitly as future work.
It is listed here as explicit future work, not as an existing result.
\end{flag}

\end{document}
